\begin{slikaDesno}{fig/LVDT.pdf}
\PID 
На слици је приказана конструкција линеарног диференцијалног трансформатора -- сензора положаја.
 Три коаксијална идентична калема, сваки дужине $h$ састављени су од равномерно и густо намотаних по $N$ навојака танке жице, 
 постављени су на пластични носач па позиционирани као на слици. Дуж осе система постављен је немагнетски штап на коме је 
 постављен магнетски добро проводан ваљак, чији се померај мери. Магнетски ваљак задире у сензорске намотаје 
 за дужине $t_1$ и $t_2$, при чему у њима остаје део испуњен ваздухом дужина $\updelta_1$ и $\updelta_2$, респективно.
 Сопстена инудктивност погонског намотаја је $L = N^2 L_0$, док су међусобне индуктивности погонског и сензорскиг намотаја дати
 изразом
 $M_{i} = N^2 \left( \dfrac{\uppi R^2}{2\updelta_i} + \dfrac{gt_i^2}{4h_i}  \right)$, где је 
 $g = 2\upmu_0 \uppi/\ln(D/d)$ релуктанса према 
 
\end{slikaDesno}
 
\begin{enumerate}[label=(\alph*)]
\item У завистности од процепа $d \ll \ell$, изразити индуктивност која се види на прикључцима намотаног калема, $L = L(d)$. 
    Магнетско расипање занемарити.
\item Скицирати зависност индуктивности од ширине процепа у интервалу вредности $100\unit{\upmu m} < d < 3\unit{mm}$. 
\end{enumerate}